% Modernised reading — fonds Alexandre Grothendieck, shelfmark 51
% (pages 1-22, the whole folder).
%
% This is an INTERPRETATION, not a transcription. It re-reads the pages in
% today's notation and vocabulary, states the hypotheses the manuscript leaves
% implicit, and drops the critical apparatus so that the mathematics reads
% straight through. Everything the brackets and underlines carried moves into a
% footnote. It is held to being mathematically correct as it stands; where the
% manuscript is loose or mistaken, the true statement is what appears here and
% the footnote says what the page has. In any dispute of substance the
% transcription governs: whoever cites Grothendieck must cite batch-01.fr.tex
% (pp. 1-20) or batch-02.fr.tex (pp. 21-22).
%
% Pass: Opus 5.5 (claude-opus-5-5), 2026-09-23 — first pass, unchecked against the pages by a human.
% The folder was read whole, its two batches together; this is one document
% for the shelfmark. It was made on Opus 5.5 and has not been measured against
% the reference reading of folder 115 (made on Opus 5). The literature named
% in the footnotes is cited from memory and was not checked against the
% sources.
%
% Scope. Page 1 is a cover whose words (« "Bonne réduction" (Shimura,
% Koizumi) ») head the document. Every other page is mathematics. Page 21
% carries « 11 » in his hand, top right; no other page is numbered by him.
%
% Spine. Good reduction of abelian varieties over a DVR, by schematic closure
% rather than by Néron models, then in families:
%   pp. 2-5    flat closure over a DVR: V-flat quotients of F <-> quotients of
%              F_K; closed subschemes; closed subgroups (Prop. 1, Cor. 1-4).
%   pp. 6-7    the scheme M of proper flat subgroup schemes, closed in Hilb;
%              proper pieces by the valuative criterion (Prop. 2).
%   pp. 8-11   isogenies extend (Prop. 3); good reduction is an isogeny
%              invariant; the Koizumi-Shimura theorem (quotients) by
%              idempotents; the closure of an abelian subvariety (Remarque).
%   pp. 12-15  in char. p the closure can fail to be abelian (two subgroups of
%              order p of a supersingular reduction); in char. 0 it cannot;
%              the Th. of p. 14 (subgroups of A_K extend uniquely over a normal
%              base); Question; Hom; the Lemma on rational sections.
%   pp. 16-17  good reduction of A_K generated by a rational map from X_K;
%              of Alb(X_K) and (Pic^0)_red.
%   pp. 18-21  the abelian part of Pic_{X/S} in a family, by a moduli problem
%              M -> S; closes only in residue char. 0 (his own margin).
%   p. 22      abelian schemes over a regular base: projectivity, torsors,
%              H^1(S,G) in H^1(K,G), purity, and the question of purity for
%              good reduction.
%
% Notation fixed for the document. V a DVR, pi, K, k; X_K the generic fibre,
% j : X_K -> X; « adhérence schématique » \overline{Y'} for the page's
% « sous-préschéma adhérence ». H' = Ker u', H its closure on pp. 8-9 (the
% page drifts from H to Gamma; footnoted). « Multiabélien » kept as his word,
% defined once (proper smooth group scheme whose neutral component is an
% abelian scheme). « Simple » (his, for smooth) rendered « lisse ». Pic, Hilb
% without his double underline. On p. 22, G the abelian scheme and P a torsor
% under it.
%
% Substantive departures, each footnoted where it occurs:
%  p. 3    quasi-coherence of J* justified (j affine); flatness proved via
%          G inside j_*G'.
%  p. 5    the unit section of H supplied.
%  p. 7    the unit condition supplied (Hilb contains the empty subscheme);
%          the « il est connu » closedness of Z in W cited (flat proper Z).
%  p. 8    H finite: argument supplied (closure of a finite set of points).
%  p. 10   the extension of p', q' to A is not justified on the page (Weil,
%          which p. 8 says is not used for Prop. 3); the margin's direct
%          route via A/C given as an alternative.
%  p. 11   « C'est un schéma abélien » read as conditional: false in char. p
%          by the example that follows; true statement: closure abelian iff
%          B -> A is a closed immersion.
%  p. 13   the illegible predicate of M_1 read as « étale ».
%  p. 14   « normal » kept (used by the Lemma); the Question answered for
%          flat A (A^0), open otherwise.
%  p. 15   the Lemma needs « géométriquement unibranche »; counterexample
%          ours (x^2 + y^2 = x^3 over R).
%  p. 16   « Pic A se réduit bien » read as A_K; hypotheses of the generic
%          curve made explicit; an étale extension of V made harmless.
%  p. 17   X_K geometrically normal supplied for (Pic^0)_red to be abelian;
%          « isogène » kept, duality suffices.
%  pp. 18-21  the Théorème as the page states it (with its hypotheses
%          struck) is false: a smooth cubic degenerating to a nodal one;
%          stated here for f smooth projective, S a reduced Q-scheme, the
%          frame in which every step of the argument holds and which the
%          page's own margins point to. Where S reduced enters, said.
%  p. 20   the last « Pic(X_K) » read as Pic of the special fibre.
%  p. 22   « K_D » read as K (or its completion at D); d) proved when the
%          order of P is invertible on S, left open otherwise.
%
% Announced and never established: the « il est connu » of p. 7; the
% illegible margins of pp. 12 and 14; the NB of p. 14 (proper flat suffices);
% the Question of p. 14 for non-flat A; the Remarque on Hom (p. 15) beyond
% the statement; the passage « il faut regarder de plus près » (p. 21), i.e.
% the Théorème in char. p; b) and d) of p. 22; the purity question.
\documentclass[11pt,a4paper]{article}
% Le préambule commun à toutes les transcriptions du fonds Grothendieck.
%
% Il tient en une idée : **l'incertitude se compose, elle ne se résout pas.**
% Une transcription qui lisse un mot douteux a détruit la seule chose pour
% laquelle on la faisait — un lecteur doit pouvoir savoir, à chaque endroit,
% ce qui était sur le feuillet et ce qui a été ajouté. Les six macros qui
% suivent sont donc l'appareil critique, et non de la décoration.
%
% Les mêmes commandes sont reconnues par `scripts/render.mjs`, qui produit la
% vue de lecture du site. Une macro ajoutée ici doit l'être aussi là-bas,
% faute de quoi le rendu échouera bruyamment — ce qui est voulu : un rendu
% silencieusement faux est pire qu'un rendu absent.

\ProvidesPackage{grothendieck}[2026/08/08 Transcriptions du fonds Grothendieck]

\usepackage{amsmath,amssymb,amsthm}
\usepackage{mathrsfs}
\usepackage{stmaryrd}
% Les diagrammes commutatifs sont partout dans ce fonds : c'est la langue
% même de la géométrie algébrique de Grothendieck, et une transcription qui
% les rendrait en prose ne transcrirait rien.
\usepackage{tikz-cd}
% Français par défaut — c'est la langue des feuillets — mais l'anglais est
% chargé aussi : la traduction et le résumé sont des documents anglais, et la
% typographie française (espaces avant les deux-points) y serait une faute.
% Ces documents basculent par \selectlanguage{english} en tête de corps.
\usepackage[english,french]{babel}
\usepackage{fontspec}
\usepackage{geometry}
\geometry{a4paper,margin=25mm,marginparwidth=28mm}
\usepackage{marginnote}
\usepackage{xcolor}
% ulem, not soul. `soul`'s \st and \ul are built on a character-by-character
% scan that XeLaTeX's font machinery breaks: the document fails with an
% undefined control sequence rather than degrading. ulem does the same two jobs
% and is Unicode-engine safe. `normalem` stops it hijacking \emph, which the
% transcriptions use for Grothendieck's own emphasis.
\usepackage[normalem]{ulem}
\usepackage{hyperref}
\hypersetup{colorlinks=true,linkcolor=black,urlcolor=black,pdfborder={0 0 0}}

% --- Métadonnées du lot -----------------------------------------------------
% Reprises telles quelles par le script de rendu, qui les affiche en tête de
% la vue de lecture. Elles fixent ce que le fichier prétend couvrir : sans
% elles, une transcription flotte sans point d'attache dans un fonds de seize
% mille feuillets.

\newcommand{\folder}[1]{\gdef\grothFolder{#1}}
\newcommand{\batch}[1]{\gdef\grothBatch{#1}}
\newcommand{\leaves}[2]{\gdef\grothFirst{#1}\gdef\grothLast{#2}}
\newcommand{\foldertitle}[1]{\gdef\grothFoldertitle{#1}}
\gdef\grothFolder{?}\gdef\grothBatch{?}\gdef\grothFirst{?}\gdef\grothLast{?}\gdef\grothFoldertitle{}

% --- Le résumé --------------------------------------------------------------
%
% Il ouvre la lecture modernisée, sous le titre « Résumé », et il est composé
% en petit. Ce n'est pas un ornement : c'est le seul passage du document qui
% ne soit pas des mathématiques, et un lecteur doit pouvoir le reconnaître
% avant de l'avoir lu — puis le sauter s'il connaît déjà le sujet, ou s'y
% arrêter s'il ne le connaît pas. Le filet qui le suit marque la reprise du
% corps.

\newenvironment{resume}
  {\small\setlength{\parskip}{0.45em}}
  {\par\medskip\hrule\medskip}

% --- Mots-clés ---------------------------------------------------------------
%
% En anglais, en clôture du résumé de la lecture modernisée : le vocabulaire
% moderne sous lequel chercher ce que les feuillets construisent. Le manifeste
% du site (`npm run manifest`) les extrait pour étiqueter la cote entière —
% cette ligne est la source unique des tags, il n'y a pas de fichier à côté.
\newcommand{\keywords}[1]{\par\smallskip\noindent
  {\footnotesize\textsc{Keywords} --- \textit{#1}}\par}

% --- L'appareil critique ----------------------------------------------------

% Le numéro de feuillet, en marge. C'est l'ancre qui permet de régler un
% désaccord sur une formule en regardant *un* feuillet plutôt qu'une chemise
% de six cents. Le site s'en sert aussi pour tourner les pages du fac-similé
% quand on fait défiler la transcription.
\newcommand{\leaf}[1]{%
  \marginnote{\footnotesize\color{gray!70}\textsf{#1}}%
  \label{leaf:#1}\ignorespaces}

% Illisible : on ne devine pas. Un mot inventé qui se lit comme les autres est
% le pire résultat possible.
\newcommand{\ill}{\textcolor{red!60!black}{[\,\ldots\,]}}

% Lecture douteuse : le mot est donné, et signalé comme incertain.
\newcommand{\uncertain}[1]{\uline{#1}}

% Ajout de l'éditeur — une lettre manquante, un mot sous-entendu.
\newcommand{\add}[1]{\textcolor{blue!50!black}{[#1]}}

% Biffé par Grothendieck lui-même. Ce qu'il a rayé fait partie du document :
% une rature dit souvent où la pensée a changé de direction.
\newcommand{\struck}[1]{\sout{#1}}

% Note du transcripteur, sur l'état matériel du feuillet.
\newcommand{\note}[1]{\par\smallskip{\small\color{gray!60!black}\textit{[#1]}}\par\smallskip}

% Note marginale de Grothendieck, distincte du corps du texte.
\newcommand{\marginal}[1]{\par\smallskip\noindent\hspace{1em}%
  {\small\color{gray!60!black}#1}\par\smallskip}

% --- Titre ------------------------------------------------------------------

\AtBeginDocument{%
  \begin{center}
    {\large\bfseries Fonds Alexandre Grothendieck --- cote n\textsuperscript{o}~\grothFolder}\\[2pt]
    {\itshape\grothFoldertitle}\\[4pt]
    {\small feuillets \grothFirst--\grothLast\ (lot \grothBatch)}\\[2pt]
    {\footnotesize\color{gray!60!black}Transcription établie d'après le fac-similé numérisé
     par l'Université de Montpellier.\\
     Les lectures incertaines sont soulignées, les illisibles notées [\,\ldots\,],
     les ajouts entre crochets.}
  \end{center}
  \vspace{1em}}

\endinput


\folder{51}
\pages{1}{22}
\dating{s.d. — le groupe « Variétés abéliennes » (45 à 56) est daté 1961-1973}
\watermark{Édition de démonstration\\interprétation personnelle de l'œuvre}
\foldertitle{« Bonne réduction » (Shimura-Koizumi) : notes manuscrites (s.d.) — lecture modernisée du dossier entier}

\begin{document}

\section*{Résumé}

\begin{resume}
Une courbe elliptique, ou plus généralement une variété abélienne, peut être
donnée par des équations à coefficients rationnels. On peut alors réduire ces
équations modulo un nombre premier $p$ et regarder ce qu'il en reste. Parfois
l'objet réduit est encore une variété abélienne de même dimension : on dit que
la variété a \emph{bonne réduction} en $p$. Parfois il dégénère, comme un
cercle qui s'écrase sur une croix quand on fait tendre un paramètre vers zéro.
Ces feuillets posent une question simple : quelles opérations respectent la
bonne réduction ? Si une variété a bonne réduction, en va-t-il de même de ses
quotients, de ses sous-variétés, des variétés qui lui sont reliées par une
isogénie, c'est-à-dire par un homomorphisme surjectif à noyau fini ? La
chemise du dossier porte deux noms, Shimura et Koizumi, qui avaient répondu
oui pour les quotients.

L'outil de Grothendieck est l'\emph{adhérence schématique}. On a un objet
au-dessus du point « générique » --- sur le corps des rationnels, disons --- et
une famille qui le contient. Son adhérence dans la famille est la plus petite
façon de le prolonger jusqu'à la réduction modulo $p$ : son ombre portée sur la
fibre spéciale. Elle existe toujours, elle est unique, et elle garde la
structure de groupe. Mais elle peut être moins bonne que l'objet de départ. En
caractéristique $p$, deux sous-groupes distincts d'ordre $p$ d'une courbe
elliptique peuvent se confondre à la réduction, comme deux points qui se
rejoignent à la limite. L'adhérence d'une sous-variété abélienne peut alors
acquérir, sur la fibre spéciale, une structure infinitésimale qui l'empêche
d'être abélienne. En caractéristique nulle cela n'arrive jamais, et le dossier
en tire un énoncé propre : sur une base normale, tout sous-groupe de la fibre
générique se prolonge d'une seule manière.

Le reste du dossier passe aux familles. D'abord la variété de Picard et celle
d'Albanese d'une variété projective, dont la bonne réduction se ramène, par
une courbe tracée sur la variété et sa jacobienne, au théorème de Shimura et
Koizumi. Ensuite une question en famille : la partie abélienne du schéma de
Picard d'une famille de variétés forme-t-elle elle-même une famille ? La
réponse passe par un problème de modules dont chaque fibre n'a qu'un point.
Grothendieck la mène à bout en caractéristique nulle, et le note en marge ;
en caractéristique $p$ il écrit qu'« il faut regarder de plus près ». La
dernière page change de cadre. Au-dessus d'une base régulière, un objet défini
génériquement ne peut refuser de se prolonger qu'en codimension un. La page se
demande alors si la bonne réduction, vérifiée le long de chaque hypersurface,
suffit partout. C'est ce qu'on appelle aujourd'hui la question de
\emph{pureté} pour les schémas abéliens.

Le dossier n'est pas daté. Son vocabulaire --- « préschéma », « simple » pour
lisse, un renvoi à un « Ch. II » pour les critères de propreté --- est celui des
premières années des \emph{Éléments}. C'est une indication, pas une date.

\keywords{good reduction, abelian scheme, schematic closure, isogeny, Picard scheme, torsor, purity}
\end{resume}

\section*{Le fil du dossier, et les conventions}

\begin{enumerate}
\item[1.] \emph{L'adhérence schématique au-dessus d'un trait} (pages 2 à 5) :
quotients plats, sous-schémas fermés plats, sous-groupes fermés plats.
\item[2.] \emph{Le schéma des sous-groupes} (pages 6 et 7), fermé dans le
schéma de Hilbert, et propre par morceaux grâce au point 1.
\item[3.] \emph{Isogénies et théorème de Koizumi--Shimura} (pages 8 à 11) : les
isogénies se prolongent, la bonne réduction est un invariant d'isogénie, les
quotients et les sous-variétés d'une variété à bonne réduction ont bonne
réduction.
\item[4.] \emph{Caractéristique $p$ contre caractéristique nulle} (pages 12 à
15) : un contre-exemple, puis le prolongement unique des sous-groupes en
caractéristique nulle, une question, une remarque sur les $\mathrm{Hom}$, et un
lemme sur les sections rationnelles.
\item[5.] \emph{Picard et Albanese} (pages 16 et 17).
\item[6.] \emph{La partie abélienne de $\mathrm{Pic}_{X/S}$} (pages 18 à 21),
démontrée en caractéristique nulle seulement.
\item[7.] \emph{Schémas abéliens sur une base régulière} (page 22).
\end{enumerate}

\emph{Conventions.} $V$ est un anneau de valuation discrète, $\pi$ une
uniformisante, $K$ son corps des fractions et $k$ son corps résiduel. Pour un
$V$-schéma $X$, on note $X_K$ la fibre générique, $j \colon X_K \to X$
l'immersion ouverte, et $X_k$ la fibre spéciale ; l'indice $K$ ou le prime
($F' = F_K$) désignent indifféremment la restriction à $X_K$, comme sur les
pages. Un \emph{schéma abélien} sur $S$ est un $S$-schéma en groupes propre et
lisse à fibres géométriquement connexes ; on dit que $A_K$ \emph{a bonne
réduction} sur $V$ s'il est la fibre générique d'un schéma abélien sur $V$ ---
la page dit qu'il « se réduit bien ». Le mot \emph{multiabélien}, qui est le
sien et n'est pas usuel, désigne ici un $S$-schéma en groupes propre et lisse
dont la composante neutre est un schéma abélien.\footnote{Le mot apparaît aux
pages 14, 15 et 21 sans définition. C'est la seule lecture compatible avec
l'usage qu'il en fait : l'adhérence d'un sous-groupe quelconque, fini ou non,
d'une variété abélienne en caractéristique nulle. Les pages écrivent aussi
« simple » au sens de « lisse », selon l'usage de l'époque ; on écrit
« lisse ».} On écrit $\mathrm{Hilb}$ et $\mathrm{Pic}$ sans le double
soulignement qu'il leur donne.

\emph{Ce que le dossier annonce et n'établit pas} : le fait « connu » sur
lequel repose la proposition 2 ; l'énoncé de la page 14 sans l'hypothèse de
normalité, et sa note marginale, illisible ; le théorème sur les
$\mathrm{Hom}$ esquissé page 15 ; le théorème des pages 18 à 21 en
caractéristique $p$ ; les affirmations b) et d) de la page 22, et la question
de pureté qui la clôt.

\pagerange{2}{5}
\section*{L'adhérence schématique au-dessus d'un trait (pages 2 à 5)}

\textbf{Proposition 1.} Soient $X$ un $V$-schéma et $F$ un $\mathcal{O}_X$-module
quasi-cohérent, $F_K = j^{*}F$.
\begin{enumerate}
\item[(i)] Soit $G' = F_K/J'$ un quotient quasi-cohérent de $F_K$. Soit $J^{*}$
le plus grand sous-module de $F$ dont la restriction à $X_K$ est contenue dans
$J'$, c'est-à-dire le noyau de $F \to j_{*}G'$, et $G = F/J^{*}$. Alors $G$ est
quasi-cohérent, plat sur $V$, et $G|_{X_K} = G'$. C'est le plus petit quotient
de $F$ qui induise sur $X_K$ un quotient majorant $G'$.
\item[(ii)] Réciproquement, si $G = F/J$ est un quotient quasi-cohérent de $F$
plat sur $V$, alors $J = J^{*}$ pour $J' = J|_{X_K}$ : $G$ est le plus petit
quotient de $F$ qui induise $G|_{X_K}$.
\end{enumerate}

\emph{Démonstration.} L'immersion $j$ est affine --- localement,
$X_K = \mathrm{Spec}\, R[1/\pi]$ --- donc $j_{*}G'$ est quasi-cohérent, et
$J^{*}$ aussi, comme noyau d'un morphisme de modules
quasi-cohérents.\footnote{La page dit seulement « on sait déjà (bof) que $J$
est q.\ coh. » ; la justification est la nôtre.} Le module $G$ s'identifie à
l'image de $F$ dans $j_{*}G'$, sur lequel $\pi$ est inversible ; $\pi$ n'est
donc pas diviseur de zéro dans $G$, ce qui, sur un anneau de valuation
discrète, est la platitude. Comme $j^{*}$ est exact et $j^{*}j_{*}G' = G'$, on
a $G|_{X_K} = G'$.\footnote{La page raisonne autrement, par minimalité : le
noyau $\mathcal{N}$ de la multiplication par $\pi$ dans $G$ est nul sur $X_K$,
donc son image inverse dans $F$ est contenue dans $J^{*}$, donc
$\mathcal{N} = 0$. Les deux arguments sont équivalents.}

Pour (ii), on a $J \subset J^{*}$ puisque $J|_{X_K} = J'$. Le quotient
$N = J^{*}/J$ est un sous-module de $G$, donc sans $\pi$-torsion, et il est nul
sur $X_K$ : toute section locale de $N$ est tuée par une puissance de $\pi$.
Donc $N = 0$.\footnote{Deux mots interlinéaires de ce passage ne se lisent pas
sûrement ; l'argument, lui, est clair sur la page : « il suffit de prouver
qu'un tel Module q.\ coh.\ est Nul. Or c'est immédiat, puisque
$N \otimes_V K = 0$ ».}

\textbf{Corollaire 1.} L'application $G \mapsto G|_{X_K}$ est une bijection de
l'ensemble des quotients quasi-cohérents de $F$ plats sur $V$ sur l'ensemble
des quotients quasi-cohérents de $F_K$.

\textbf{Corollaire 2.} Soient $X_1$, $X_2$ deux $V$-schémas et $F_1$, $F_2$
des modules quasi-cohérents sur eux. La bijection du corollaire 1 est
compatible au produit tensoriel externe : si $G_i$ est le quotient plat de
$F_i$ qui correspond à $G'_i$, le quotient plat de $F_1 \boxtimes F_2$ sur
$X_1 \times_V X_2$ qui correspond à $G'_1 \boxtimes G'_2$ est
$G_1 \boxtimes G_2$.\footnote{La page écrit $F_1 \otimes_V F_2$, pour le
produit tensoriel externe sur $X_1 \times_V X_2$.} En effet $G_1 \boxtimes G_2$
est un quotient de $F_1 \boxtimes F_2$, plat sur $V$ comme produit tensoriel
de modules plats, de restriction $G'_1 \boxtimes G'_2$ ; on applique (ii).

\textbf{Corollaire 3.} Les sous-schémas fermés $Y$ de $X$ plats sur $V$
correspondent bijectivement aux sous-schémas fermés $Y'$ de $X_K$, par
$Y \mapsto Y_K$ et $Y' \mapsto \overline{Y'}$, l'\emph{adhérence
schématique} de $Y'$ dans $X$ : le plus petit sous-schéma fermé de $X$ dont la
restriction à $X_K$ contient $Y'$. Cette correspondance est compatible aux
produits fibrés sur $V$ :
$\overline{Y'_1 \times_K Y'_2} = \overline{Y'_1} \times_V \overline{Y'_2}$.

C'est le cas $F = \mathcal{O}_X$ des corollaires 1 et 2. L'espace sous-jacent
à $\overline{Y'}$ est l'adhérence de $Y'$ dans $X$, d'où le
nom.\footnote{C'est aujourd'hui EGA~IV$_{2}$, §~2.8 (cité de mémoire), où
l'adhérence schématique et sa platitude au-dessus d'un trait sont traitées
dans ces termes. La page écrit « sous-préschéma adhérence ».}

\textbf{Corollaire 4.} Soient $G$ un $V$-schéma en groupes et $H'$ un
sous-schéma en groupes fermé de $G_K$. Alors $H = \overline{H'}$ est un
sous-schéma en groupes fermé de $G$, plat sur $V$, de fibre générique $H'$, et
c'est le seul.

\emph{Démonstration.} L'unicité est celle du corollaire 3. L'inversion de $G$
est un automorphisme qui préserve $H'$, donc son adhérence. Pour la
multiplication $m \colon G \times_V G \to G$, le sous-schéma fermé
$m^{-1}(H)$ contient $H' \times_K H'$ sur la fibre générique, donc contient son
adhérence, qui est $H \times_V H$ par le corollaire 3. Enfin la section unité
$e \colon \mathrm{Spec}\, V \to G$ se factorise par $H$ : $e^{-1}(H)$ est un
sous-schéma fermé de $\mathrm{Spec}\, V$ qui contient le point générique, donc
c'est $\mathrm{Spec}\, V$.\footnote{La page traite l'inversion (« évident ») et
la multiplication, et omet la section unité ; l'argument pour celle-ci est le
nôtre.}

\pagerange{6}{7}
\section*{Le schéma des sous-groupes (pages 6 et 7)}

\textbf{Proposition 2.} Soient $S$ un schéma localement noethérien et $G$ un
$S$-schéma en groupes dont le schéma de Hilbert $\mathrm{Hilb}_{G/S}$ existe ;
il paramètre les sous-schémas fermés de $G$ propres et plats sur
$S$.\footnote{Une note marginale demande « $G$ q.-projectif sur $S$ ??? » :
c'est l'hypothèse sous laquelle on sait construire ce schéma.} Pour tout
$S' \to S$, soit $\mathcal{M}(S')$ l'ensemble des sous-schémas fermés de
$G_{S'}$, propres et plats sur $S'$, qui sont des sous-schémas en groupes.
\begin{enumerate}
\item[(a)] Le sous-foncteur $\mathcal{M}$ de $\mathrm{Hilb}_{G/S}$ est
représentable par un sous-schéma fermé de $\mathrm{Hilb}_{G/S}$.
\item[(b)] Si $G$ est propre sur $S$ et si $\mathrm{Hilb}_{G/S}$ est réunion
disjointe d'ouverts $H_i$ de type fini sur $S$, alors $\mathcal{M}$ est réunion
disjointe des $\mathcal{M}_i = \mathcal{M} \cap H_i$, et chaque
$\mathcal{M}_i$ est propre sur $S$.
\end{enumerate}

\emph{Démonstration de (a).} Soit $H \subset G_{T}$ le sous-schéma universel
au-dessus de $T = \mathrm{Hilb}_{G/S}$. Il faut montrer que les $T' \to T$ au-dessus
desquels $H_{T'}$ est un sous-groupe sont ceux qui se factorisent par un
sous-schéma fermé de $T$. Être un sous-groupe se dit par trois inclusions :
la section unité $e$ se factorise par $H$ ; $H$ est contenu dans son image
$\check{H}$ par l'inversion ; $H \times_T H$ est contenu dans
$P = m^{-1}(H)$. Or si $Z$ est un sous-schéma fermé propre et plat sur la base
et $W$ un sous-schéma fermé quelconque, les points de la base au-dessus
desquels $Z \subset W$ forment un sous-schéma fermé.\footnote{C'est le fait
que la page dit « connu » : l'inclusion $Z \subset W$ équivaut à la nullité
d'un morphisme de modules cohérents à valeurs dans $\mathcal{O}_Z$, qui est
plat et propre sur la base, et le lieu où un tel morphisme s'annule est fermé
(EGA~III$_{2}$, §~7.7, cité de mémoire).} Pour la section unité, c'est
$e^{-1}(H)$. On prend pour $\mathcal{M}$ l'intersection des trois sous-schémas
fermés ainsi obtenus.\footnote{La page n'exprime que les deux dernières
conditions et prend $T = T_1 \cap T_2$. La première est nécessaire : le
schéma de Hilbert contient le sous-schéma vide, qui est stable par inversion
et par multiplication sans être un sous-groupe. Sur une fibre non vide, les
deux autres l'entraînent.}

\emph{Démonstration de (b).} Chaque $H_i$, ouvert d'une réunion disjointe, est
aussi fermé ; $\mathcal{M}_i$ est donc fermé dans $H_i$, de type fini et
séparé sur $S$. Reste le critère valuatif d'existence. Si
$\mathrm{Spec}\, V \to S$ est un trait et $H'$ un sous-groupe fermé propre de
$G_K$, son adhérence schématique dans $G_V$ est un sous-groupe fermé plat
(corollaire 4 de la proposition 1), propre puisque $G$ l'est : c'est un
$V$-point de $\mathcal{M}$ qui prolonge le $K$-point donné, et il reste dans
$\mathcal{M}_i$ puisque $\mathrm{Spec}\, V$ est connexe. C'est pour cela
que la proposition 1 est là.\footnote{La page renvoie en une ligne à « la
caractérisation habituelle des propriétés » et au « Cor.~4 de prop.~1. OK. » ;
on a déplié.}

\emph{Variantes.} La page note encore que les sous-groupes finis, et de même
les sous-groupes abéliens, forment des parties ouvertes de $\mathcal{M}$, la
première étant aussi fermée.\footnote{Une troisième classe, entre les deux, est
lue « séparables » sans certitude. Ces affirmations ne sont pas démontrées sur
la page. Pour les sous-groupes finis, elles tiennent parce que, pour un
schéma en groupes plat et de présentation finie, la dimension des fibres est
localement constante (SGA~3, exposé VI$_{B}$, cité de mémoire). Pour les
sous-groupes abéliens, la lissité est une condition ouverte, et la connexité
des fibres est ouverte et fermée dans une famille propre et lisse.}

\pagerange{8}{11}
\section*{Isogénies, et le théorème de Koizumi--Shimura (pages 8 à 11)}

\textbf{Proposition 3.} Soient $A$ un schéma abélien sur $V$ et
$u' \colon A_K \to B'$ une isogénie vers une variété abélienne $B'$ sur $K$.
Il existe un schéma abélien $B$ sur $V$, avec $B_K \simeq B'$, et un
homomorphisme $u \colon A \to B$ qui induit $u'$.

La page souligne que la démonstration n'utilise pas le théorème de prolongement
de Weil, selon lequel un $K$-morphisme $A_K \to B_K$ entre fibres génériques de
schémas abéliens provient d'un morphisme $A \to B$.

\emph{Démonstration.} Le noyau $H' = \operatorname{Ker} u'$ est un sous-groupe
fini de $A_K$. Son adhérence $H = \overline{H'}$ est un sous-groupe fermé de
$A$, plat sur $V$ (corollaire 4), propre, et quasi-fini. En effet chaque point
de $H'$ a un corps résiduel fini sur $K$, et son adhérence dans $H$ est de
dimension au plus $1$ ; elle ne rencontre donc la fibre spéciale qu'en un
nombre fini de points. Ainsi $H$ est fini et plat sur $V$.\footnote{La page
dit seulement « propre et quasi-fini sur $V$ : fibre générale finie ».
L'argument de dimension qui donne la finitude de la fibre spéciale est le
nôtre. La page appelle ce sous-groupe $H$, $H'$ dans les deux premières lignes,
puis $\Gamma$, $\Gamma'$ à partir de « $A/\Gamma$ ». Les deux glyphes se
ressemblent, et c'est vraisemblablement la même lettre ; on écrit $H$
partout.} Comme $A$ est projectif sur $V$, le quotient $B = A/H$ existe. C'est
un $V$-schéma en groupes, plat et propre, et $A \to B$ a pour noyau
$H$.\footnote{La justification de la projectivité est illisible sur la page.
Le quotient d'un schéma projectif par un groupe fini et plat qui opère
librement existe (SGA~3, exposé V, cité de mémoire).} La formation de ce
quotient commute au changement de base. Donc $B_K = A_K/H' \simeq B'$ et
$B_k = A_k/H_k$, qui est une variété abélienne. Un schéma en groupes plat et
propre dont les fibres sont des variétés abéliennes est un schéma
abélien.\footnote{La page écrit « $A \otimes k / \Gamma' \otimes k$ », avec le
prime ; c'est $H_k$, la fibre spéciale de l'adhérence.}

\textbf{Corollaire 1.} Si $A'$ et $B'$ sont deux variétés abéliennes isogènes
sur $K$, $A'$ a bonne réduction si et seulement si $B'$ a bonne réduction, et
alors les deux réductions sont isogènes.

En effet, toute isogénie $u'$ admet une isogénie $v'$ en sens inverse (avec
$v' u' = n$ pour un entier $n$), et on applique la proposition 3 dans les deux
sens. Le prolongement $u \colon A \to B$ a un noyau fini $H$, et sa réduction
$u_k$ est surjective par égalité des dimensions : c'est une isogénie.

\emph{Conséquence.} La bonne réduction est une propriété de la classe
d'isogénie. Un produit de variétés à bonne réduction a bonne réduction. En
particulier, si les facteurs simples de $A'$ --- au sens du théorème de
complète réductibilité de Poincaré --- ont bonne réduction, $A'$ a bonne
réduction. La réciproque est le théorème suivant.

\textbf{Théorème (Koizumi--Shimura).} Si $A'$ a bonne réduction sur $V$, tout
quotient abélien de $A'$ a bonne réduction, et de même toute sous-variété
abélienne de $A'$.\footnote{Shimura et Koizumi ont établi cet énoncé dans le
langage des spécialisations de Weil (G.~Shimura et S.~Koizumi, articles de 1959
sur la spécialisation des variétés abéliennes, cités de mémoire). Dans le
langage actuel, c'est un corollaire du critère de Néron--Ogg--Šafarevič : la
bonne réduction se lit sur la non-ramification du module de Tate, qui passe
aux sous-quotients (Serre--Tate, 1968). La démonstration de ces pages n'utilise
ni modèle de Néron ni module de Tate, seulement l'adhérence schématique et des
idempotents. La page énonce le théorème pour les quotients ; les
sous-variétés viennent de la remarque de la page 11.}

\emph{Démonstration.} Soit $B'$ un quotient abélien de $A'$. Par la complète
réductibilité de Poincaré, $A'$ est isogène à $B' \times C'$ pour une variété
abélienne $C'$, et par le corollaire 1 on peut supposer
$A' = B' \times C'$. On peut donc supposer que le quotient est un facteur
direct. Il y a alors des endomorphismes $p'$, $q'$ de $A'$ tels que
\[
  p' + q' = \mathrm{id}, \qquad p'^{2} = p', \qquad q'^{2} = q'.
\]
Ils se prolongent en endomorphismes $p$, $q$ de $A$, qui satisfont les mêmes
identités, puisque $\mathrm{End}(A) \to \mathrm{End}(A_K)$ est
bijectif.\footnote{La page écrit « les endomorphismes correspondants $p$, $q$
de $A$ » sans dire pourquoi ils existent. C'est le théorème de prolongement de
Weil --- celui-là même que le N.B.\ de la page 8 dit ne pas utiliser pour la
proposition 3. La note marginale de la page 10, « on peut le prouver
directement comme proposition 3 », indique l'autre route, qui s'en passe.
Soit $C'$ le noyau de $A' \to B'$ et $C$ son adhérence, un sous-groupe fermé
plat de $A$. Le quotient $A/C$ est plat et propre sur $V$, de fibres les
quotients de variétés abéliennes $A_K/C'$ et $A_k/C_k$, donc abélien ; c'est le
prolongement cherché. Il faut seulement savoir que le quotient d'un schéma
abélien par un sous-groupe fermé plat existe au-dessus d'un trait
(Anantharaman, 1973, cité de mémoire).} On en déduit un isomorphisme de
schémas en groupes
\[
  A \xrightarrow{\ \sim\ } \operatorname{Ker} p \times_V \operatorname{Ker} q
  = A_1 \times_V A_2, \qquad x \mapsto (q x,\ p x),
\]
d'inverse $(y, z) \mapsto y + z$. Chaque $A_i$ est un rétracte de $A$, donc
lisse et propre sur $V$. Comme $A_k = (A_1)_k \times (A_2)_k$ est connexe, les
deux facteurs le sont. Ce sont des schémas abéliens, et $B'$ est la fibre
générique de l'un d'eux.

\textbf{Remarque (page 11).} Soient $A$ un schéma abélien sur $V$ et $B'$ une
sous-variété abélienne de $A_K$. D'après ce qui précède, $B' = B_K$ pour un
schéma abélien $B$ sur $V$, et l'inclusion $B' \to A_K$ se prolonge en un
homomorphisme $v \colon B \to A$, par le théorème de Weil. L'image
schématique de $v$ est l'adhérence $\overline{B'}$. En effet l'image induit
$B'$, donc contient $\overline{B'}$. Inversement $v^{-1}(\overline{B'})$ est un
sous-schéma fermé de $B$ qui contient $B_K$, donc c'est $B$ puisque $B$ est
plat. Ainsi $\overline{B'}$ est un sous-groupe fermé de $A$, plat sur $V$.
Les conditions suivantes sont alors équivalentes : $\overline{B'}$ est un
schéma abélien ; $B \to \overline{B'}$ est un isomorphisme ; $v$ est une
immersion fermée ; le noyau de $v$ est trivial.\footnote{La page écrit
« C'est un schéma abélien, et $B \to B'$ est un isom. (grâce au th.\ de Weil
p.ex.), i.e.\ $B \to A$ est un monomorphisme ». Lu comme une affirmation, c'est
faux, et la page le dit elle-même à la ligne suivante : « D'après
Shimura--Koizumi, il n'en est pas toujours ainsi ». On lit donc la phrase comme
l'énoncé des équivalences. Sur la page, « $B \to B'$ » est écrit sans barre ;
l'argument demande $B \to \overline{B'}$.}

\pagerange{12}{15}
\section*{Caractéristique $p$, caractéristique nulle (pages 12 à 15)}

\subsection*{Le contre-exemple}

Ces conditions peuvent être en défaut. Supposons $k$ de caractéristique $p > 0$
et $K$ de caractéristique nulle. Soit $B$ une courbe elliptique sur $V$ dont la
réduction $B_k$ n'a pas de point d'ordre $p$, c'est-à-dire est
supersingulière. Alors $B_k$ n'a qu'un sous-groupe d'ordre $p$, le noyau du
Frobenius. Quitte à remplacer $K$ par une extension finie, $B_K$ a deux
sous-groupes distincts $\Gamma'_1$, $\Gamma'_2$ d'ordre $p$. Leurs adhérences
$\Gamma_1$, $\Gamma_2$ sont finies et plates d'ordre $p$, et leurs fibres
spéciales coïncident. Ainsi $\Gamma'_1 \cap \Gamma'_2 = (e)$ mais
$\Gamma_1 \cap \Gamma_2 \neq (e)$. On pose
\[
  A = B/\Gamma_1 \times_V B/\Gamma_2
\]
et on y envoie $B$ diagonalement. Sur la fibre générique c'est une immersion
fermée, de noyau $\Gamma'_1 \cap \Gamma'_2$ ; sur la fibre spéciale, non. Par
la remarque précédente, l'adhérence de $B_K$ dans $A$ est un sous-groupe plat
qui n'est pas un schéma abélien : sa fibre spéciale n'est pas
réduite.\footnote{La page conclut par « OK » après la construction ; la
conséquence sur l'adhérence est celle de la remarque qu'elle illustre. Une
note marginale en oblique affirme que l'exemple marche encore pour $B$ abélien
de dimension $\geq 2$ pourvu que la puissance $p$-ième soit nulle dans
$\mathrm{Lie}(B)$. Elle est trop illisible pour être reconstituée, et on ne le
fait pas. Trois lignes lourdement biffées précèdent la construction.}

\subsection*{La caractéristique nulle}

En caractéristique nulle, tout schéma en groupes de type fini sur un corps est
lisse (théorème de Cartier). L'adhérence d'un sous-groupe fermé de $A_K$, plate
à fibres lisses, est donc lisse. Si le sous-groupe est une sous-variété
abélienne, les fibres géométriques de l'adhérence restent connexes, puisque
leur nombre de composantes est localement constant dans une famille propre et
lisse. L'adhérence est alors un schéma abélien. Plus généralement, l'adhérence
d'un sous-groupe quelconque de $A_K$ est multiabélienne.\footnote{La page dit
« tout sous-schéma en groupes fermé $B'$ de $A_K$ est la restriction d'un
ss-schéma abélien $B$ de $A$ » : « abélien » ne vaut que si $B'$ est connexe,
et l'énoncé de la page 14 dit « multiabélien ». L'argument sur les
composantes connexes est le nôtre.}

Soit maintenant $S$ localement noethérien de caractéristique nulle et $A$ un
schéma abélien projectif sur $S$, de sorte que $\mathrm{Hilb}_{A/S}$ existe et
est localement de type fini. Le schéma $\mathcal{M}$ de la proposition 2 est
alors localement propre sur $S$ : c'est la proposition 2~(b), l'adhérence d'un
sous-groupe restant, en caractéristique nulle, un sous-groupe lisse. De plus
$\mathcal{M}$ est non ramifié sur $S$. Un sous-groupe fini est contenu dans un
$A[n]$, qui est fini étale, et ses sous-groupes forment un schéma fini étale.
Une sous-variété abélienne ne se déforme pas comme sous-groupe.\footnote{Pour
le schéma $\mathcal{M}_{1}$ des sous-groupes finis, l'attribut est illisible
sur la page. L'interligne « loc.\ non ramifié » le suggère, et c'est vrai :
$\mathcal{M}_1$ est même étale. La rigidité des sous-variétés abéliennes est
énoncée en caractéristique nulle par les pages 13 et 14, sans démonstration.}

\textbf{Théorème (page 14).} Soient $S$ un schéma localement noethérien,
connexe et normal, de caractéristique nulle, de corps des fonctions $K$, et
$A$ un schéma abélien sur $S$ tel que $\mathrm{Hilb}_{A/S}$ existe et soit
localement de type fini (par exemple, $A$ projectif sur $S$). Pour tout
sous-schéma en groupes fermé $B'$ de $A_K$, il existe un unique sous-schéma en
groupes fermé multiabélien $B$ de $A$ tel que $B_K = B'$. Si $B'$ est
abélien, $B$ l'est.

\emph{Démonstration.} Les composantes $\mathcal{M}_i$ de $\mathcal{M}$ sont
propres et non ramifiées sur $S$, donc finies et non ramifiées. Le sous-groupe
$B'$ est une section rationnelle de l'une d'elles. Par le lemme ci-dessous,
elle est partout définie, et la section obtenue $S \to \mathcal{M}_i$ est le
$B$ cherché. L'unicité vient de ce que $\mathcal{M}$ est séparé et $S$
réduit.\footnote{La page met un point d'interrogation après « normal », et se
demande dans un alinéa très raturé si « $S$ réduit » suffirait (« probable,
mais… »). La normalité sert au lemme. On laisse la question où la page la
laisse. Une note marginale, presque illisible, affirme qu'il suffit que $A$
soit propre et plat sur $S$, à composante neutre multiabélienne.}

\textbf{Question (page 14).} En caractéristique nulle, un $S$-schéma en groupes
propre $A$ a-t-il un plus grand sous-schéma abélien ?\footnote{Si $A$ est plat
sur $S$, oui. Il est alors lisse (fibres lisses par Cartier), sa composante
neutre $A^{0}$ est un sous-groupe ouvert à fibres connexes (SGA~3,
exposé VI$_{B}$, cité de mémoire), propre, donc un schéma abélien. Tout
sous-schéma abélien, à fibres connexes contenant l'unité, est contenu dans
$A^{0}$. Sans platitude, c'est la vraie question, et c'est exactement celle que
posent les pages 18 à 21 pour la partie propre de $\mathrm{Pic}_{X/S}$.}

\textbf{Remarque (page 15).} Pour $A$, $B$ multiabéliens sur $S$, la situation
est analogue pour $\mathrm{Hom}_{S\text{-gr}}(A, B)$, et cette fois sans
hypothèse de caractéristique. Ce schéma, s'il existe et est localement de type
fini, est non ramifié et localement propre sur $S$. D'où un théorème de
prolongement des homomorphismes analogue au précédent.\footnote{Aujourd'hui,
pour des schémas abéliens sur une base normale, l'application
$\mathrm{Hom}(A, B) \to \mathrm{Hom}(A_K, B_K)$ est bijective (Faltings--Chai,
chap.~I, cité de mémoire). C'est le théorème que la page annonce, et qu'elle
n'écrit pas.}

\textbf{Lemme (page 15).} Soit $X$ un schéma localement noethérien,
géométriquement unibranche, et soit $f \colon X' \to X$ un morphisme fini et non
ramifié. Toute section rationnelle de $f$ est partout définie.

\emph{Démonstration.} On se ramène à $X = \mathrm{Spec}\, R$ intègre et à $X'$
l'adhérence de l'image de la section. Alors $X'$ est intègre, et
$f \colon X' \to X$ est fini, non ramifié et birationnel, défini par une
$R$-algèbre $R'$ contenue dans la clôture intégrale de $R$ dans son corps des
fractions. Comme $X$ est géométriquement unibranche, les fibres de $f$ ont un
seul point géométrique. Un morphisme fini, non ramifié et radiciel est une
immersion fermée. Birationnelle, elle est un isomorphisme.\footnote{La page
écrit « unibranche », après avoir biffé « normal ». L'hypothèse géométrique est
nécessaire. La cubique $x^{2} + y^{2} = x^{3}$ sur $\mathbb{R}$ est unibranche
à l'origine, dont les deux branches sont conjuguées sur $\mathbb{C}$. Sa
normalisation est finie, non ramifiée et birationnelle, mais la section
rationnelle inverse ne se prolonge pas en l'origine. Pour le théorème de la
page 14, « normal » suffit, puisqu'un schéma normal est géométriquement
unibranche. Tout le lemme est entre crochets sur la page, et deux mots de la
fin de la démonstration sont illisibles.}

\pagerange{16}{17}
\section*{Picard et Albanese (pages 16 et 17)}

\textbf{Proposition (page 16).} Soient $A_K$ une variété abélienne sur $K$, $X$
un schéma projectif et plat sur $V$, à fibre générique géométriquement
intègre, dont les fibres sont lisses en dehors d'une partie de codimension au
moins $2$. Soit $u \colon X_K \dashrightarrow A_K$ une application rationnelle
qui engendre $A_K$, au sens où les différences $u(x) - u(y)$ engendrent $A_K$.
Alors $A_K$ a bonne réduction sur $V$.

\emph{Démonstration.} Quitte à faire une extension non ramifiée de $V$, ce qui
ne change pas la question,\footnote{Par le critère de Néron--Ogg--Šafarevič, la
bonne réduction ne change pas par extension non ramifiée. La page écrit
seulement « sur un corps alg.\ clos ».} on choisit une « courbe générique »
$C \subset X$, intersection complète d'hyperplans généraux. Elle évite le lieu
singulier de $X_k$, de codimension au moins $2$ dans $X_k$, et le lieu
d'indétermination de $u$. $C$ est donc une courbe lisse et propre sur $V$, à
fibres géométriquement connexes, et $C_K \to A_K$ est un morphisme qui engendre
encore $A_K$.\footnote{La page dit « on sait que cela existe… » et, pour la
seconde propriété, « je dis qu'on sait que ». Les deux sont des énoncés de
type Bertini et Lefschetz qu'elle ne démontre pas. Le mot lu « successivement »
dans « successivement simple » est une conjecture sur un mot réduit à un
trait, et la lettre lue $X$ dans « dans $X$ » ressemble à un $V$.} Par la
propriété universelle de la jacobienne, on obtient un homomorphisme surjectif
$\mathrm{Pic}^{0}(C_K) \to A_K$. Or $\mathrm{Pic}^{0}(C_K)$ a bonne réduction :
c'est la fibre générique du schéma abélien $\mathrm{Pic}^{0}_{C/V}$. Par le
théorème de Koizumi--Shimura, $A_K$ a bonne réduction.\footnote{La page conclut
« $\mathrm{Pic}\, A$ se réduit bien ». C'est $A_K$ qui est en cause, et
l'énoncé l'exige. Un bloc biffé, encadré, tentait la route duale, par
$\mathrm{Pic}^{0}(A_K) \to \mathrm{Pic}^{0}(C_K)$.}

\textbf{Théorème (page 17).} Soit $X$ un schéma projectif et plat sur $V$, dont
les fibres sont lisses en dehors d'une partie de codimension au moins $2$, et
dont la fibre générique est géométriquement intègre et géométriquement
normale. Alors la variété d'Albanese $\mathrm{Alb}(X_K)$ et la variété de
Picard $(\mathrm{Pic}^{0}_{X_K/K})_{\mathrm{red}}$ ont bonne réduction sur
$V$.\footnote{L'hypothèse de normalité géométrique est ajoutée : c'est elle qui
fait de $(\mathrm{Pic}^{0})_{\mathrm{red}}$ une variété abélienne. La page ne
suppose que les fibres « simples en codim $\leq 1$ » et $X_K$ géométriquement
irréductible. L'exposant de $\mathrm{Pic}^{0}(X_K)$ est surchargé ; on le lit
« red ».}

\emph{Démonstration.} La page rappelle d'abord que, s'il existe dans
$\mathrm{Pic}^{0}(X_K)$ un sous-schéma abélien de même ensemble sous-jacent,
$\mathrm{Alb}(X_K)$ est défini comme son dual.\footnote{« S'il existe » n'est
pas une précaution de style. Sur un corps imparfait, le sous-schéma réduit
d'un schéma en groupes n'est pas toujours un sous-groupe.} Quitte à faire une
extension non ramifiée de $V$, $X$ a une section, qui passe par un point lisse
de $X_k$ (lemme de Hensel). Cette section donne un point rationnel de $X_K$,
et avec lui le morphisme d'Albanese $X_K \to \mathrm{Alb}(X_K)$, qui engendre.
Par la proposition précédente, $\mathrm{Alb}(X_K)$ a bonne réduction, et son
dual $(\mathrm{Pic}^{0})_{\mathrm{red}}$ aussi : le dual d'un schéma abélien
est un schéma abélien.\footnote{La page conclut par l'isogénie entre les deux
(« comme $\mathrm{Pic}^{0}(X_K)'$ est isogène »), ce qui est vrai aussi, via
une polarisation, et suffit par le corollaire 1 de la proposition 3.}

\pagerange{18}{21}
\section*{La partie abélienne du schéma de Picard (pages 18 à 21)}

La page énonce : \emph{si $f \colon X \to S$ est projectif, $S$ localement
noethérien et réduit, il existe un unique sous-schéma abélien $A$ de
$\mathrm{Pic}_{X/S}$ dont l'ensemble sous-jacent soit la réunion $P$ des
composantes neutres des fibres de $\mathrm{Pic}_{X/S}$.} Les hypothèses
« plat, à fibres géométriques intègres et simples en codimension $1$ » ont été
écrites puis biffées. Deux notes marginales restreignent l'énoncé :
« Prouvé seulement en car.~$0$ pour l'instant… » et, au bas de la page 20,
« Attention, cela est-il vrai si $f \colon X \to S$ n'est pas simple
?! ».\footnote{Tel qu'il est écrit, l'énoncé est faux. Une famille de cubiques
planes lisses qui dégénère en une cubique nodale est projective, plate, à
fibres géométriquement intègres, au-dessus d'une courbe lisse. La composante
neutre du Picard de la fibre nodale est $\mathbf{G}_m$, qui n'est pas propre, et
$P$ n'est la réunion des fibres d'aucun schéma abélien. Même les hypothèses
biffées n'y suffiraient pas : il faut au moins des fibres géométriquement
normales pour que les $\mathrm{Pic}^{0}$ réduits soient propres, et la page le
sent (« comme ses fibres sont propres sur $k$ [vrai …] »).}

On énonce donc le théorème dans le cadre où chaque étape de l'argument
tient, et que les marges désignent.

\textbf{Théorème.} Soit $f \colon X \to S$ un morphisme projectif et lisse à
fibres géométriquement connexes, $S$ étant un schéma localement noethérien
réduit de caractéristique nulle. Il existe un unique sous-schéma abélien $A$
de $\mathrm{Pic}_{X/S}$ dont l'ensemble sous-jacent est la réunion $P$ des
composantes neutres des fibres.\footnote{Dans ce cadre, on sait aujourd'hui
que $\mathrm{Pic}^{0}_{X/S}$ est lui-même un schéma abélien, grâce à la
théorie de Hodge, qui rend $R^{1}f_{*}\mathcal{O}_X$ localement libre et
compatible au changement de base (Kleiman, \emph{The Picard scheme}, §~9.5,
cité de mémoire) ; $A = \mathrm{Pic}^{0}_{X/S}$, et l'hypothèse « $S$ réduit »
devient inutile. La seconde note marginale de la page 18 dit déjà que cette
hypothèse n'est pas nécessaire, en renvoyant à « la rigidité des Modules
formels des V.A. ». Le cas de la caractéristique $p$, que la page 21 laisse
ouvert, n'est pas traité ici.}

\emph{Démonstration.} C'est un problème de modules. Les fibres de $P$ sont
propres, donc $P$ est propre sur $S$.\footnote{« par un théorème connu
[renvoyant aux critères de propreté de Ch.~II] » : sans doute le chapitre~II
des \emph{Éléments}. La page se demande en même temps s'il en est de même de
$\mathrm{Pic}^{\tau}_{X/S}$, et si $\mathrm{Pic}_{X/S}$ est réunion disjointe
d'ouverts propres. Une marge répond « ? Non ».} Pour $S' \to S$, soit
$\mathcal{F}(S')$ l'ensemble des sous-schémas fermés de $P_{S'}$, plats sur
$S'$, qui sont des sous-groupes de $\mathrm{Pic}_{X_{S'}/S'}$, abéliens, et
ensemblistement égaux à $P_{S'}$. Ce foncteur est représentable par un
$S$-schéma $\mathcal{M}$. Sans la dernière condition, c'est un sous-schéma de
$\mathrm{Hilb}_{P/S}$ (proposition 2), où la condition « abélien » est
ouverte. La condition ensembliste est elle aussi ouverte : elle est stable par
changement de base, et, vraie en un point, elle l'est au voisinage par raison
de dimension.

(i) \emph{Les fibres géométriques de $\mathcal{M}$ sont réduites à un point
réduit.} Sur un corps algébriquement clos, la seule sous-variété abélienne de
$\mathrm{Pic}$ qui ait pour ensemble sous-jacent $(\mathrm{Pic}^{0})_{\mathrm{red}}$
est celle-ci, et une sous-variété abélienne est rigide. La page invoque « le
th.\ de rigidité ». Donc $\mathcal{M} \to S$ est radiciel et non ramifié :
c'est un monomorphisme.

(ii) \emph{$\mathcal{M} \to S$ est propre.} Il est de type fini et séparé. Reste
le critère valuatif. Soit $S = \mathrm{Spec}\, V$ un trait, et donnons-nous
une section rationnelle de $\mathcal{M}$, c'est-à-dire la sous-variété
abélienne $A_K$ de $\mathrm{Pic}^{0}(X_K)$ de même ensemble sous-jacent. Par le
théorème de la page 17, $A_K$ a bonne réduction : c'est la fibre générique
d'un schéma abélien $A$ sur $V$.\footnote{La page 20 appelle ce résultat
« le théorème de Koizumi ». C'est le théorème de la page 17, ou l'énoncé de
Koizumi sur la spécialisation des variétés de Picard et d'Albanese dont il
procède ; on ne l'a pas vérifié.} L'inclusion $A_K \to \mathrm{Pic}_{X_K}$ se
prolonge en un morphisme $u \colon A \to P$.\footnote{La page invoque ici « le
théorème de connexion de Murre--Chow », qu'elle n'énonce pas, et ajoute : « On
n'a pas besoin du théorème si $X$ simple sur $S$ », ce qui est notre cadre.}
L'image de $u$ est fermée, de dimension relative $\dim A_K$. Or
$\dim A_K = \dim \mathrm{Pic}^{0}(X_K) = \dim \mathrm{Pic}^{0}(X_k)$, par la
constance de la dimension des fibres du Picard, que la lissité de $f$
garantit.\footnote{La page écrit au dernier terme $\dim \mathrm{Pic}(X_K)$ ;
c'est la fibre spéciale qui est en cause. La marge « Attention, cela est-il vrai
si $f \colon X \to S$ n'est pas simple ?! » vise précisément ce pas.} L'image
de $u$ est donc $P$ (page 21).

Reste à voir que le noyau $N$ de $u$, fini sur $S$, est trivial. Sa fibre
spéciale $N_0$ est un $p$-groupe, $p$ étant l'exposant caractéristique de
$k$. En effet, pour $n$ premier à $p$, $\mathrm{Pic}_{X/S}[n]$ est non ramifié
sur $S$ et $A[n]$ est un revêtement étale de $S$. Le noyau de
$A[n] \to \mathrm{Pic}[n]$ est donc ouvert et fermé dans $A[n]$, de fibre
générique triviale, donc trivial. Si $k$ est de caractéristique nulle, $p = 1$
et « on a fini » : $u$ est un monomorphisme propre, donc une immersion fermée
d'image $P$, et c'est la section cherchée.\footnote{La page note $K$ le noyau
et $K_0$ sa fibre spéciale. On écrit $N$ pour ne pas le confondre avec le corps
$K$. Un N.B.\ encadré, annulé par deux diagonales, passait par une section
d'ordre $n$ relevée à $S$. La page 21 porte en haut à droite le chiffre
« 11 », de sa main, le seul numéro d'auteur du dossier ; on ne sait pas ce
qu'il compte.} Sinon, écrit-il, « il faut regarder de plus près ».

(iii) \emph{Conclusion.} Un monomorphisme propre est une immersion fermée. Elle
est surjective, puisqu'au-dessus de chaque point de $S$ --- de corps résiduel
parfait en caractéristique nulle --- $(\mathrm{Pic}^{0})_{\mathrm{red}}$ est une
sous-variété abélienne. Comme $S$ est réduit, $\mathcal{M} = S$, et l'objet
universel au-dessus de $\mathcal{M}$ est le $A$ cherché. C'est là, et là
seulement, que sert l'hypothèse « $S$ réduit ».

Sous un trait horizontal, la page ouvre un programme : étudier noyaux et
images d'homomorphismes de schémas abéliens et multiabéliens en
caractéristique $\neq 0$, notamment sur une base non réduite.\footnote{Le signe
« $\neq$ » surcharge une autre lettre et reste douteux ; un point
d'interrogation est dans la marge.}

\pagerange{22}{22}
\section*{Schémas abéliens sur une base régulière (page 22)}

Le titre est de sa main. $S$ est un schéma régulier intègre, de corps des
fonctions $K$. Pour un diviseur premier $D$ de $S$, $\mathcal{O}_D$ est
l'anneau local de $S$ au point générique de $D$, un anneau de valuation
discrète.

a) \emph{Un schéma abélien sur $S$ est projectif sur $S$}, par un « raisonnement
de Nakai- » dont le second nom est illisible.\footnote{La page écrit « Un
schéma abélien $P$ », le mot « abélien » barré d'un grand trait vertical. La
suite traite $P$ comme un torseur, c'est-à-dire un schéma abélien sans section
unité choisie. Que les schémas abéliens sur une base normale soient projectifs
localement sur la base est un théorème de Raynaud (\emph{Faisceaux amples sur
les schémas en groupes}, 1970, cité de mémoire). Sur une base régulière,
l'argument naturel prolonge un diviseur ample de la fibre générique par
adhérence, ce qui est possible parce que les diviseurs de Weil y sont de
Cartier, et contrôle l'amplitude fibre à fibre par un critère numérique de type
Nakai. C'est vraisemblablement le raisonnement visé, mais c'est une
reconstruction. Pour un torseur, nous ne savons pas dire si l'énoncé tient en
général.}

b) Pour un tel $P$, il existe un schéma abélien $G$ « de translations », et $P$
est un torseur isotrivial sous $G$.\footnote{Affirmé sans démonstration. Le
groupe $H^{1}(S, G)$ est de torsion sur une base normale, puisqu'il s'injecte
dans $H^{1}(K, G)$ (point c). Un torseur d'ordre $n$ provient d'un torseur sous
$G[n]$, qui se trivialise sur un revêtement fini, étale si $n$ est inversible
sur $S$. Au-delà, nous ne le démontrons pas.}

c) Si $G$ est un schéma abélien sur $S$, l'application
\[
  H^{1}(S, G) \longrightarrow H^{1}(K, G)
\]
est injective. Un torseur $P_K$ sous $G_K$ provient d'un torseur sous $G$ si et
seulement s'il « se réduit bien ». Il y a un plus grand ouvert dense $U$ de $S$
au-dessus duquel $P_K$ se prolonge.\footnote{L'injectivité vaut dès que $S$ est
normal. Un torseur qui a une section rationnelle a une section définie en
codimension $1$, par le critère valuatif de propreté, puis partout, par le
théorème de prolongement de Weil pour les applications rationnelles vers un
schéma en groupes lisse (Bosch--Lütkebohmert--Raynaud, \emph{Néron Models},
§~4.4, cité de mémoire). Par la même injectivité, les prolongements locaux se
recollent, d'où le plus grand ouvert $U$. Le mot lu « ouvert » a sur la page la
forme d'un trait terminé en flèche.}

d) Le fermé $S - U$ est de codimension $1$ : c'est ce que la marge rattache au
« th.\ de pureté ».\footnote{C'est-à-dire : si $P_K$ se prolonge au-dessus d'un
ouvert dont le complémentaire est de codimension $\geq 2$, il se prolonge
partout. Quand l'ordre $n$ de $P_K$ est inversible sur $S$, c'est le théorème de
pureté de Zariski--Nagata appliqué au torseur sous $G[n]$, qui est fini étale
(SGA~1, exposé X, cité de mémoire). Quand $n$ n'est pas inversible, nous ne
savons pas dire si l'énoncé tient. La fin de la phrase de la page,
« codim.~$1$ en … », est illisible.}

Il en résulte une injection
\[
  H^{1}(K, G) / H^{1}(S, G) \;\hookrightarrow\;
  \bigoplus_{D} H^{1}(K, G) / H^{1}(\mathcal{O}_D, G),
\]
la somme portant sur les diviseurs premiers de $S$. La somme est directe parce
qu'un torseur sur $K$ se prolonge au-dessus d'un ouvert dense, qui ne manque
qu'un nombre fini de diviseurs. L'application est injective : une classe qui
se prolonge en chaque $D$ a un $U$ qui contient tous les points de codimension
$1$, et d) force $U = S$.\footnote{La page écrit $H^{1}(K_D, G)$, l'indice lu
sans certitude. Comme $\mathcal{O}_D \subset K$, le corps des fractions de
$\mathcal{O}_D$ est $K$ lui-même. Une colonne de mots soulignés, à laquelle
mène une flèche partie de ce terme, se lit « étudier … complétion ». Elle
suggère le passage au complété de $\mathcal{O}_D$, que la page n'écrit pas.
L'inclusion est notée « $\subset$ » sur la page, et $\sum$ pour la somme.}

La page s'achève sur une question, prise dans une boucle : existe-t-il un
analogue du théorème de pureté pour la bonne réduction des schémas abéliens ?
Autrement dit, si $A_K$ a bonne réduction en chaque diviseur premier d'une base
régulière, a-t-il bonne réduction sur la base entière ?\footnote{La question a
sa littérature propre, sous le nom de « pureté » des schémas abéliens. Quand les
caractéristiques résiduelles ne divisent pas l'entier en jeu, le critère de
Néron--Ogg--Šafarevič et la pureté de Zariski--Nagata prolongent le module de
Tate. Le prolongement du schéma abélien lui-même est traité notamment dans
Faltings--Chai (1990, chap.~V) et par Vasiu et Zink (2010), références citées
de mémoire, dont nous ne résumons pas les énoncés.}

\end{document}
